\pdfvariable trailerid {[<C3372026090300000000000000000000><C3372026090300000000000000000000>]}
\documentclass[10pt]{article}
\usepackage[margin=0.72in]{geometry}
\usepackage{amsmath,amssymb,amsthm,mathtools,microtype,hyperref}
\hypersetup{colorlinks=true,linkcolor=blue,urlcolor=blue}
\providecommand{\CRevisionRound}{2}
\newtheorem{theorem}{Theorem}
\newtheorem{lemma}[theorem]{Lemma}
\newtheorem{remark}[theorem]{Remark}
\newcommand{\HH}{\mathcal H}
\newcommand{\ZZ}{\mathbb Z}
\newcommand{\EE}{\mathbb E}
\newcommand{\opn}{\widehat n}
\newcommand{\ii}{\mathrm i}
\setlength{\emergencystretch}{3em}
\title{Integer Quantum Resonance and Antiresonance of the Kicked Rotor:\\
An Exact Parity Sheet}
\author{Route-A source-local certificate HCS-C337}
\date{3 September 2026}

\begin{document}
\maketitle
\begin{abstract}
We solve the positive-integer resonance sheet of the quantum kicked rotor in
a fixed free-after-kick convention.  Parity of the integer resonance index
gives a complete dichotomy: an exact Bessel propagator with ballistic
momentum variance, or a state-independent two-kick antiresonance.
\ifnum\CRevisionRound>0
On the resonant face we also derive the full momentum characteristic
function, survival probability, and centered moments through order six.
\fi
\ifnum\CRevisionRound>1
All momentum seeds, zero kick strength, zero elapsed kicks, vector phases,
operator ordering, source ownership, and Route-A exclusions are closed
explicitly; general rational resonance and detuning remain outside scope.
\fi
\end{abstract}

\section{Frozen Floquet owner}
Let
\[
 \HH=L^2(\mathbb T,d\theta/(2\pi)),\qquad
 |n\rangle(\theta)=e^{\ii n\theta},\qquad
 \opn=-\ii\partial_\theta.
\]
The periodic $H^1$ realization of $\opn$ is self-adjoint, and $\opn^2$
is self-adjoint on periodic $H^2$.  For $\kappa\in\mathbb R$ put
\[
 K_\kappa=e^{-\ii\kappa\cos\theta},\qquad
 U_\tau=e^{-\ii\tau\opn^2/2}K_\kappa.
\]
Here $K_\kappa$ denotes unit-modulus multiplication.  Spectral calculus and
bounded multiplication make both factors, hence $U_\tau$, unitary.  The
ordering is part of the definition: one kick followed by free rotation.

We restrict to
\begin{equation}\label{eq:sheet}
 \tau=2\pi\ell,\qquad \ell\in\ZZ_{>0}.
\end{equation}
No assertion below concerns another rational sheet or a detuned value.

\begin{lemma}[Free-phase parity]\label{lem:parity}
At~\eqref{eq:sheet}, the free factor is $I$ for even $\ell$ and the half-turn
\[
 (Rf)(\theta)=f(\theta+\pi)
\]
for odd $\ell$.
\end{lemma}

\begin{proof}
On $|n\rangle$ the free eigenvalue is $e^{-\ii\pi\ell n^2}$.  Since
$n^2-n=n(n-1)$ is even, this equals $e^{-\ii\pi\ell n}$.  It is one for
even $\ell$ and $(-1)^n$ for odd $\ell$.  The latter is precisely the
Fourier action of $R$.
\end{proof}

\begin{theorem}[Integer resonance and antiresonance]\label{thm:main}
Fix $m\in\ZZ$ and $t\in\ZZ_{\geq0}$.
\begin{enumerate}
\item If $\ell$ is even, then
\begin{equation}\label{eq:kernel}
 \langle n|U_{2\pi\ell}^{t}|m\rangle
 =(-\ii)^{n-m}J_{n-m}(\kappa t).
\end{equation}
Consequently the momentum law is $J_{n-m}(\kappa t)^2$, its mean is $m$,
and
\begin{equation}\label{eq:ballistic}
 \operatorname{Var}(n)=\frac{\kappa^2t^2}{2},\qquad
 \EE_m[\opn^2/2]_t
 =\frac{m^2}{2}+\frac{\kappa^2t^2}{4}.
\end{equation}
\item If $\ell$ is odd, then $U_{2\pi\ell}=RK_\kappa$ and
\begin{equation}\label{eq:involution}
 U_{2\pi\ell}^{2}=I.
\end{equation}
At even $t$ the amplitude is $\delta_{nm}$, while at odd $t$ it is
\begin{equation}\label{eq:odd-kernel}
 (-1)^n(-\ii)^{n-m}J_{n-m}(\kappa).
\end{equation}
\end{enumerate}
\end{theorem}

\begin{proof}
For even $\ell$, Lemma~\ref{lem:parity} gives $U=K_\kappa$ and
$U^t=K_{t\kappa}$.  The Jacobi--Anger expansion
\[
 e^{-\ii x\cos\theta}
 =\sum_{q\in\ZZ}(-\ii)^qJ_q(x)e^{\ii q\theta}
\]
with $q=n-m$ proves~\eqref{eq:kernel}.  The moment statements follow from
Fourier orthogonality:
\[
 \sum_q J_q(x)^2e^{\ii qu}=J_0\!(2x\sin(u/2)).
\]
Two derivatives at $u=0$ give mean displacement zero and variance $x^2/2$;
translation by $m$ gives~\eqref{eq:ballistic}.

For odd $\ell$, the free factor is $R$.  Since the half-turn reverses the
cosine,
\[
 RK_\kappa R=K_{-\kappa}=K_\kappa^{-1}.
\]
Thus $(RK_\kappa)^2=I$.  Even powers are $I$, and odd powers are $RK_\kappa$.
Applying $R$ after the Jacobi--Anger expansion multiplies the $n$th Fourier
coefficient by $(-1)^n$, which proves~\eqref{eq:odd-kernel}.
\end{proof}

The phrase \emph{operator parity owner} records the first paper round: the
all-time dichotomy is an operator theorem, not a finite-lattice simulation.

\ifnum\CRevisionRound>0
\section{Exact momentum statistics}
Let $x=\kappa t$ on the even sheet and set $q=n-m$.  If
$c_q=(-\ii)^qJ_q(x)$, Fourier orthogonality gives
\begin{align}
 \Phi_x(u)
 &=\sum_{q\in\ZZ}|c_q|^2e^{\ii qu}\label{eq:char}\\
 &=\frac1{2\pi}\int_0^{2\pi}
 e^{-\ii x\cos(\theta+u)}e^{\ii x\cos\theta}\,d\theta
 =J_0\!(2x\sin(u/2)).\nonumber
\end{align}
The last equality follows by shifting the integration variable in the
standard integral representation of $J_0$; its sign is immaterial because
$J_0$ is even.  At $u=0$,~\eqref{eq:char} also proves normalization.

\begin{theorem}[Characteristic and moment closure]\label{thm:moments}
For even $\ell$, the survival probability is
\[
 \Pr_m\{n_t=m\}=J_0(\kappa t)^2,
\]
and the centered odd moments through order six vanish.  The nonzero centered
moments are
\begin{align}
 \mu_2&=\frac{x^2}{2},\label{eq:moments}\\
 \mu_4&=\frac{x^2}{2}+\frac{3x^4}{8},\nonumber\\
 \mu_6&=\frac{x^2}{2}+\frac{15x^4}{8}+\frac{5x^6}{16}.\nonumber
\end{align}
For odd $\ell$, the law is the point mass at $m$ at even times and the
one-kick Bessel law $J_{n-m}(\kappa)^2$ at odd times; use $x=0$ and
$x=\kappa$, respectively, in~\eqref{eq:moments}.
\end{theorem}

\begin{proof}
The survival formula is~\eqref{eq:kernel} at $n=m$.  For a characteristic
function, $\EE(q^r)=\Phi_x^{(r)}(0)/\ii^r$.  Six direct differentiations of
the final expression in~\eqref{eq:char} give~\eqref{eq:moments}; evenness in
$u$ annihilates the odd derivatives.  The odd-sheet statement follows from
Theorem~\ref{thm:main}.
\end{proof}

The phrase \emph{Bessel characteristic owner} marks the substantive first
revision.  It upgrades a variance statement to an exact distributional
identity and moments through sixth order.
\fi

\ifnum\CRevisionRound>1
\section{Boundary atlas and convention audit}
At $t=0$, $J_q(0)=\delta_{q0}$ closes both parity formulas.  At $\kappa=0$,
the even sheet has $U=I$, while the odd sheet has $U=R$.  Hence every
momentum probability is stationary.  On the odd sheet a vector $|m\rangle$
acquires $(-1)^m$ after one kick, although its ray is fixed; this distinguishes
vector equality from measurement equality.  For nonzero $\kappa$, the
even-sheet variance in~\eqref{eq:ballistic} is exactly quadratic at every
integer time.  Negative and zero momentum seeds need no separate argument,
because every formula is translation-covariant in $m$.

Equation~\eqref{eq:involution} is an operator identity on all of $\HH$.
It does not say that every vector has least period two: eigenvectors of the
involution may return projectively after one step.  Conversely, we do not
replace $e^{-\ii\tau\opn^2/2}K_\kappa$ by the reversed product.  Such a
convention can be related by a strobe change, but its raw amplitudes are not
silently asserted here.

\section{Evidence and collision audit}
The certificate contains 396 exact parity rows, 435 Gaussian-rational
formal Fourier coefficients, 882 exact moment rows, 120 operator-word rows,
and seven 90-digit Bessel spot checks.  It audits 22,444 scalar leaves.  A
producer-independent checker performs 47,531 assertions; an independent
SymPy lane proves 13,188 identities; two isolated generations are byte
identical; and 133 repaired-hash, parser, theorem, boundary, and evaluator
attacks are rejected.  High-precision Bessel sums are labeled
\texttt{NUMERICAL\_OBSERVATION}; the infinite identities follow from the
Fourier proof, not from truncation.

The closest workspace owners are C110 (classical nonautonomous H\'enon
Floquet dynamics), C143 (coined quantum walk), C148 (open Walsh quantum
baker), C178 (harmonic metaplectic strobe), C224 (two-level Landau--Zener
scattering), C318 (static SSH lattice), and C323 (finite continuous-time
quantum search).  None owns the cosine-kicked rotor's infinite momentum
lattice and its exact integer parity sheet.  This is a workspace collision
statement, never a claim of literature priority.

\section{Source boundary and Route-A decision}
Dana, Eisenberg and Shnerb identify quantum antiresonance with exact periodic
recurrence~\cite{Dana1996}.  Dana and Dorofeev study general kicked-particle
resonances and exact transport quantities~\cite{DanaDorofeev2006}; Ryu et al.
experimentally report rational-Talbot-time resonance and ballistic momentum
transfer~\cite{Ryu2006}.  These primary records establish provenance and
physical context.  The displayed theorem is independently derived here, and
no priority claim is made for its ingredients or assembly.

The phrase \emph{parity boundary and route firewall} marks the final
revision.  Under \texttt{NO\_BAD\_EULER\_OR\_ROOT\_NUMBER}, the strict tuple is
\[
 (\mathrm{A0\_FAIL},\mathrm{A1\_FAIL},\mathrm{A2\_FAIL},
  \mathrm{A3\_FAIL},\mathrm{A4\_NATURAL\_QUANTIZATION}).
\]
The Floquet unitary is a natural source quantization with the physical kick
clock, but its recurrences are not an isolated arithmetic primitive-orbit
ledger.  There is no prime-power clock, orbit Euler product, target Fredholm
determinant, target divisor, functional equation, Weil compression,
target-zero match, or Hilbert--Polya operator.  Route A is rejected and Route
B remains locked.  No target arithmetic local data, Euler factor, root
number, or automorphy object is asserted.

\paragraph{AI use.}
A generative language model assisted proof organization, code scaffolding,
and manuscript drafting.  The analytic derivation, producer-independent
checker, symbolic lane, hostile tests, and deterministic artifacts define
the audit record.
\fi

\begingroup
\small
\begin{thebibliography}{9}
\bibitem{Dana1996}
I. Dana, E. Eisenberg, and N. Shnerb,
``Antiresonance and localization in quantum dynamics,''
\emph{Physical Review E} 54 (1996), 5948--5963.
DOI: \href{https://doi.org/10.1103/PhysRevE.54.5948}%
{10.1103/PhysRevE.54.5948}.

\bibitem{DanaDorofeev2006}
I. Dana and D. L. Dorofeev,
``General Quantum Resonances of the Kicked Particle,''
\emph{Physical Review E} 73 (2006), 026206.
DOI: \href{https://doi.org/10.1103/PhysRevE.73.026206}%
{10.1103/PhysRevE.73.026206};
\href{https://arxiv.org/abs/nlin/0509035}{arXiv:nlin/0509035}.

\bibitem{Ryu2006}
C. Ryu, M. F. Andersen, A. Vaziri, M. B. d'Arcy, J. M. Grossman,
K. Helmerson, and W. D. Phillips,
``High-Order Quantum Resonances Observed in a Periodically Kicked
Bose--Einstein Condensate,''
\emph{Physical Review Letters} 96 (2006), 160403.
DOI: \href{https://doi.org/10.1103/PhysRevLett.96.160403}%
{10.1103/PhysRevLett.96.160403}.
\end{thebibliography}
\endgroup
\end{document}
